\section{Related Work}
\label{sec:related}

\subsection{Geospatial Foundation Models in Remote Sensing}
Pre-trained foundation models have revolutionized Earth observation by establishing general-purpose visual representations that transfer across diverse downstream tasks~\cite{zhang2024rs5m,simeoni2025dinov3,congo2023satmae}. In contrast to traditional supervised deep neural networks trained on narrow, task-specific datasets, foundation models are trained on massive archives of satellite imagery using self-supervised learning (SSL) or contrastive vision-language pre-training (VLP)~\cite{radford2021learning}.

In vision-language modeling, RemoteCLIP~\cite{wang2023remoteclip}, GeoRSCLIP~\cite{zhang2024rs5m}, and SkyCLIP~\cite{mai2023skyclip} align satellite multi-spectral imagery with descriptive domain-specific text captions derived from crowdsourced geospatial databases (e.g., OpenStreetMap) and automated vision-language bootstrapping. The fundamental contrastive objective minimizes the symmetric InfoNCE loss across a mini-batch of $B$ image-text pairs $\{(\bm{x}_i, \bm{c}_i)\}_{i=1}^B$:
\begin{equation}
\begin{split}
    \mathcal{L}_{\text{VLP}} = -\frac{1}{2B} \sum_{i=1}^B \Bigg[ \log \frac{\exp(\bm{f}_i^\top \bm{t}_i / \tau_0)}{\sum_{j=1}^B \exp(\bm{f}_i^\top \bm{t}_j / \tau_0)} \\
    + \log \frac{\exp(\bm{t}_i^\top \bm{f}_i / \tau_0)}{\sum_{j=1}^B \exp(\bm{t}_i^\top \bm{f}_j / \tau_0)} \Bigg],
\end{split}
\end{equation}
where $\bm{f}_i = \theta_{\text{opt}}(\bm{x}_i)$ and $\bm{t}_i = \theta_{\text{text}}(\bm{c}_i)$ denote normalized visual and textual embeddings, and $\tau_0$ is a learnable temperature. These models enable zero-shot classification and zero-shot open-vocabulary retrieval by computing cosine similarity between satellite image embeddings and text prompt representations.

Concurrently, masked autoencoding architectures---such as SatMAE~\cite{congo2023satmae}, Prithvi-100M/300M~\cite{jakubik2023prithvi}, Scale-MAE~\cite{reed2023scalemae}, and RingMo~\cite{qi2020mlrsnet}---reconstruct masked spatial-spectral patches from multi-spectral Sentinel-2 and Landsat pixels, capturing deep geometric and physical dependencies. Furthermore, self-supervised vision transformers such as DINOv2 and DINOv3~\cite{simeoni2025dinov3} demonstrate exceptional feature granularity, dense spatial correspondence, and out-of-distribution robustness. In this work, we harness frozen GeoRSCLIP and DINOv3 representations, completely eliminating the prohibitive cost of full-parameter fine-tuning while exploiting their rich geospatial representations.

\subsection{Multi-Source Remote Sensing for Hydrological Surveillance}
Surface water extraction and flood inundation mapping have historically relied on handcrafted multi-spectral indices, such as the Normalized Difference Water Index (NDWI)~\cite{mcfreeters1996use}, Modified NDWI (MNDWI)~\cite{xu2006modification}, and Automated Water Extraction Index (AWEI)~\cite{pekel2016high}. While optical spectral indices operate effectively under clear skies, they fail catastrophically in the presence of cloud cover, cloud shadows, and haze~\cite{rahimzadegan2021surface}.

To achieve all-weather day-and-night surface water surveillance, Synthetic Aperture Radar (SAR) systems, such as Sentinel-1 C-band radar, emit microwave pulses and record backscatter intensity ($\sigma^0_{\text{VV}}, \sigma^0_{\text{VH}}$)~\cite{clement2018multi,boudjit2021sen12flood}. Calm open water acts as a specular reflector, scattering microwave energy away from the radar antenna and appearing as dark patches in SAR backscatter imagery. However, wind-induced surface waves, capillary roughness, urban double-bounce reflections, and speckle noise frequently cause severe classification ambiguity. 

The Cloude-Pottier eigenvalue-eigenvector decomposition of the coherency matrix $\mathbf{T}_3$ represents backscattering mechanisms via entropy $H$, anisotropy $A$, and mean alpha angle $\bar{\alpha}$:
\begin{equation}
    \bar{\alpha} = \sum_{i=1}^3 P_i \alpha_i, \qquad H = -\sum_{i=1}^3 P_i \log_3 P_i,
\end{equation}
where $P_i$ are normalized eigenvalues, distinguishing surface scattering ($\bar{\alpha} \to 0^\circ$), double bounce ($\bar{\alpha} \to 90^\circ$), and volume scattering ($\bar{\alpha} \approx 45^\circ$). Recent multi-source data fusion frameworks~\cite{hong2021spectralformer,plaza2009recent} demonstrate that combining complementary optical multi-spectral and microwave SAR observations provides the most reliable foundation for water resources monitoring. However, previous fusion methods employ static, heuristic weighting that cannot adapt when one modality is degraded.

\subsection{Inductive vs. Transductive Few-Shot Adaptation}
Few-shot learning aims to recognize novel target categories given only a tiny labeled support set $\mathcal{S}$~\cite{snell2017prototypical,vinyals2016matching,sung2018learning,tian2020rethinking}. In inductive few-shot learning, the model adapts parameters solely based on the support set $\mathcal{S}$ and evaluates query samples $\bm{x} \in \mathcal{Q}$ independently. Inductive baselines include Prototypical Networks (ProtoNet)~\cite{snell2017prototypical}, Linear Probing (LP++)~\cite{huang2024lpplusplus}, and training-free key-value cache adapters such as Tip-Adapter~\cite{zhang2022tipadapter}.

In contrast, transductive few-shot learning processes the entire unlabeled query batch $\mathcal{Q}$ jointly at inference time~\cite{boudriki2020transductive,zanella2024boosting,li2026timplusplus,elkhoury2026lctim}. In remote sensing, transductive learning is particularly practical because satellite data is acquired in large geographic tiles comprising thousands of contiguous pixels or patches. Transductive Information Maximization (TIM++)~\cite{li2026timplusplus} maximizes the mutual information between query representations and predictions while penalizing divergence from zero-shot textual priors. Locally Consistent TIM (LC-TIM)~\cite{elkhoury2026lctim} augments TIM++ with a local consistency regularizer across a $k$-nearest-neighbor graph. However, LC-TIM enforces a static graph topology ($\kappa=5$) and constant regularizer weight ($\lambda_{\text{LC}}=0.3$), leading to severe label bleeding across complex hydrological boundaries.

\subsection{Reinforcement Learning in Earth Observation}
Reinforcement learning (RL) has emerged as a powerful paradigm for autonomous decision-making in remote sensing, including satellite constellation scheduling, agile Earth observation tasking, and autonomous UAV path planning~\cite{sutton2018reinforcement,schulman2017proximal,khan2022deep}. The policy gradient theorem states that the gradient of expected trajectory reward $J(\theta)$ is:
\begin{equation}
    \nabla_\theta J(\theta) = \mathbb{E}_{\pi_\theta} \left[ \nabla_\theta \log \pi_\theta(\mathbf{a} \mid \mathbf{s}) \, Q^{\pi_\theta}(\mathbf{s}, \mathbf{a}) \right].
\end{equation}
In vision-language processing, RLPrompt~\cite{sun2022rlprompt} employs policy gradients to optimize discrete textual prompts without backpropagation into the language model. Proximal Policy Optimization (PPO)~\cite{schulman2017proximal} provides stable policy gradient optimization with clipped surrogate objectives. Unlike prior works that apply RL to full model parameter fine-tuning, \rllctim{} employs RL as a lightweight, plug-and-play meta-controller for dynamic manifold graph construction and multi-source sensor gating.

\subsection{Synthesis and Distinctions of the Proposed Approach}
While previous literature has addressed either transductive adaptation or multi-modal fusion in isolation, \rllctim{} is the first framework to synthesize both mechanisms under an adaptive reinforcement-learned control policy. By treating the construction of the transductive manifold graph as a discrete Markov Decision Process, our approach eliminates heuristic hyperparameter tuning and guarantees mathematically bounded convergence over frozen foundation embeddings.
